\subsection{Research}\label{app:research}

This subsection proves both parts of Proposition~\ref{prop:research}. It also states and proves
Remark~\ref{rem:researchties} and Proposition~\ref{prop:researchregion}, the research parameter region.

% STRUCTURE: the research tuple and condition (E7) <- Section 5
\paragraph{Equilibrium with research.}
Market makers and the challenger do not observe the research choice; they use a conjectured intensity
\(\hat\lambda\). An equilibrium with research is a tuple \((\hat\lambda,\sigma_H,\sigma_L,\mu,P,\bar\mu,\zeta)\),
together with the bids. Conditions (E1)--(E6) of Definition~\ref{def:equilibrium} hold at the intensity
\(\hat\lambda\). Condition (E5) holds also at \(\hat\lambda=0\): a deviation in research makes a success possible,
and the order after it must be optimal. A seventh condition completes the definition.
\begin{enumerate}[label=(E\arabic*),start=7,leftmargin=3em]
\item \emph{Research optimality.} Hold \(P\) and \(\zeta\) at their values in the tuple. Then \(\hat\lambda\) is
the research choice of a best response among all joint deviations in research and orders.
\end{enumerate}
A joint deviation changes the true order probabilities, but not the prices or the entry rule. With the research
value \(\Psi\) below, the expected trading payoff after a success at the candidate's prices,
Lemma~\ref{lem:a5-joint} shows that the best research choice is unique and equals \(\min\{\lb,\Psi/a\}\). Condition
(E7) therefore reads \(\hat\lambda=\min\{\lb,\Psi/a\}\), and research never has a tie.

\paragraph{Standing assumptions.}
The auction is the benchmark, the prior is \(1/2\), and \(0<\beta\le1/3\). Fix a
candidate with conjectured intensity \(\hat\lambda\in[0,\lb]\) and order distributions \(\sigma_H,\sigma_L\) after a
success. Write \(\hat b=(1-\hat\lambda)\beta\) and \(\hat b_0=(1-\hat\lambda)(1-2\beta)\), and
\(f_j(q)=(1-\hat\lambda)\eta(q)+\hat\lambda\sigma_j(q)\). The conditions (E1) to (E6) hold at the intensity
\(\hat\lambda\), and (E7) is research optimality. A candidate satisfies (E1) to (E4) and (E6). Lemma~\ref{lem:trading}
applies at \(\hat\lambda\) and gives the payoffs of an order after a success,
\begin{equation}\label{eq:a5-payoffs}
U_H(q)=q\,e(q)\DT[1-\mu(q)]-k|q|,\qquad U_L(q)=-q\,e(q)\DT\mu(q)-k|q| .
\end{equation}
Let \(\Psi=\tfrac12\max_qU_H(q)+\tfrac12\max_qU_L(q)\ge0\) be the research value. Since \(r>\ell\), \(\DT>0\) and
\(g_H>g_L\) at every strength of this subsection.

\paragraph{Derived functions.}
% STRUCTURE: u_lambda, M_lambda, m_lambda and the closed form of lambda_E <- Section 5
For \(\lambda\in[0,1)\), write \(u_\lambda=\lambda/[(1-\lambda)\beta]\), \(M_\lambda=(1+u_\lambda)/(2+u_\lambda)\)
and \(m_\lambda=1-M_\lambda\); these are the posterior bounds of Section~\ref{sec:inference} at the prior \(1/2\) and
the presence \(\lambda\). Then \(m_\lambda=(1-\lambda)\beta/[2\beta+\lambda(1-2\beta)]\) falls strictly from \(1/2\)
to \(0\), because \(1-2\beta>0\). At the cap of calibration B, \(u_{\lb}=u\), so
\(M_{\lb}=\eu/(1+\eu)\) and \(m_{\lb}=1/(1+\eu)\). For \(1/2<\tau(r)<1\), the recruitment intensity
\(\lambda_E(r)\) of Section~\ref{sec:research} solves \(M_\lambda=\tau(r)\):
\[
\lambda_E(r)=\frac{\beta\,[2\tau(r)-1]}{\beta+(1-2\beta)[1-\tau(r)]},
\]
and \(d\lambda_E/d\tau=\beta/\{\beta+(1-2\beta)(1-\tau)\}^2>0\). Clear the denominator in \eqref{eq:research}:
\begin{equation}\label{eq:a5-F}
\begin{aligned}
\Lambda_r(\lambda)&=2a(1-2\beta)\lambda^2+\bigl[4a\beta+k(1-2\beta)+\beta\DT(r)\bigr]\lambda+\beta\bigl[2k-\DT(r)\bigr]\\
&=2\bigl[2\beta+\lambda(1-2\beta)\bigr]\bigl[a\lambda-\tfrac12\bigl(\DT(r)m_\lambda-k\bigr)\bigr].
\end{aligned}
\end{equation}
Every coefficient except the constant is positive, so \(\Lambda_r\) rises strictly on \(\lambda\ge0\). If \(\DT(r)>2k\),
then \(\Lambda_r(0)<0<\Lambda_r(1)\), and \(\lambda_R(r)\) is the unique root of \(\Lambda_r\) in \((0,1)\). Define the recruitment
margin and the threshold entry
\[
\varphi_a(r)=\DT(r)[1-\tau(r)]-k-2a\lambda_E(r),\qquad e^*(r)=\frac{2a\lambda_E(r)+k}{\DT(r)[1-\tau(r)]}.
\]
Since \(m_{\lambda_E}=1-\tau\), \eqref{eq:a5-F} gives \(\Lambda_r(\lambda_E(r))=-[2\beta+\lambda_E(1-2\beta)]\,\varphi_a(r)\).
Hence \(\lambda_R(r)\ge\lambda_E(r)\) exactly when \(\varphi_a(r)\ge0\), and \(e^*(r)<1\) exactly when
\(\varphi_a(r)>0\). At calibration B, \(B_r(1/2)=5-r/4+1/(8r)<c_H\) for \(1<r<10\), so \(\tau(r)>1/2\) at every
strength below.

\begin{lemma}[Joint deviations]\label{lem:a5-joint}
For every candidate and every \(z\in[0,\lb]\), the best payoff of a joint deviation \((z,\sigma'_H,\sigma'_L)\) is
\(z\Psi-(a/2)z^2\). The unique best research choice is \(z^*=\min\{\lb,\Psi/a\}\). Hence (E7) holds if and only if
\(\hat\lambda=\min\{\lb,\Psi/a\}\), and research never has a tie.
\end{lemma}

\begin{proof}
A joint deviation changes the true order law to \((1-z)\eta+z\sigma'_j\). It does not change the prices or the
preparation rule, because market makers and the challenger use the candidate's \(\hat\lambda\) and \(\sigma_j\).
The investor trades only after a success, so its expected payoff is
\[
\Pi(z,\sigma')=z\Bigl[\tfrac12\sum_q\sigma'_H(q)U_H(q)+\tfrac12\sum_q\sigma'_L(q)U_L(q)\Bigr]-\frac a2z^2 ,
\]
with \(U_j\) from \eqref{eq:a5-payoffs} at the candidate's prices and entry. The bracket is at most \(\Psi\), with
equality exactly when each \(\sigma'_j\) puts weight only on maximizers of \(U_j\). So
\(\sup_{\sigma'}\Pi(z,\sigma')=z\Psi-(a/2)z^2\). This function of \(z\) is strictly concave with derivative \(\Psi-az\). If
\(\Psi\ge a\lb\), the derivative is positive on \([0,\lb)\), and the maximizer is \(\lb\). If \(\Psi<a\lb\), the
maximizer is \(\Psi/a\in[0,\lb)\). Strict concavity makes the maximizer unique, and it also excludes a nondegenerate
mixture over research intensities.
\end{proof}

\begin{lemma}[Bounds under every profile]\label{lem:a5-bounds}
Let \(0\le\hat\lambda\le\lb<1\). For every profile and every candidate:
\begin{enumerate}[label=(\alph*),leftmargin=2em]
\item every posterior satisfies \(m_{\hat\lambda}\le\mu(q)\le M_{\hat\lambda}\), and \(M_{\hat\lambda}\le M_{\lb}\);
\item the price reveals the posterior: if \(P(q)=P(q')\), then \(\mu(q)=\mu(q')\), and \(\bar\mu(P(q))=\mu(q)\);
\item \(U_H(-1)\le-k\) and \(U_L(1)\le-k\), so every equilibrium has \(\sigma_H(-1)=\sigma_L(1)=0\);
\item if \(\sigma_H(-1)=\sigma_L(1)=0\), then \(\mu(-1)\le1/2\le\mu(1)\) and \(\mu(0)\le1/2\) when also
\(\sigma_H(1)=1\); and \(\mu(1)=M_{\hat\lambda}\) when \(\sigma_H(1)=1\).
\end{enumerate}
\end{lemma}

\begin{proof}
(a) This is Lemma~\ref{lem:beliefs}(b)--(c) at the presence \(\hat\lambda\) and the prior \(1/2\); the zero order
uses \(\beta\le1/3\). The map \(\lambda\mapsto u_\lambda\) rises, so \(M_{\hat\lambda}\le M_{\lb}\). (b) This is
Lemma~\ref{lem:prices}(c). (c) This is Lemma~\ref{lem:trading}(a). (d) Here \(f_H(1)=\hat b+\hat\lambda\sigma_H(1)\ge
\hat b=f_L(1)\) and \(f_H(-1)=\hat b\le\hat b+\hat\lambda\sigma_L(-1)=f_L(-1)\). With \(\sigma_H(1)=1\),
\(f_H(0)=\hat b_0\le\hat b_0+\hat\lambda\sigma_L(0)=f_L(0)\), and
\(\mu(1)=(\hat b+\hat\lambda)/(2\hat b+\hat\lambda)=M_{\hat\lambda}\).
\end{proof}

\begin{lemma}[Positive research needs a strict full purchase]\label{lem:a5-buy}
Let \(\tau(r)>1/2\), and take an equilibrium with \(\hat\lambda>0\). Then \(U_H(1)>0\) and \(\sigma_H(1)=1\). Also
\(\mu(1)=M_{\hat\lambda}\), \(\mu(0)\le1/2\), \(\mu(-1)\le1/2\), and \(\zeta(P(0))=\zeta(P(-1))=0\).
\end{lemma}

\begin{proof}
By Lemma~\ref{lem:a5-joint}, \(\hat\lambda=\min\{\lb,\Psi/a\}>0\), so \(\Psi>0\). By Lemma~\ref{lem:a5-bounds}(c) and
\(U_j(0)=0\), either \(U_H(1)>0\) or \(U_L(-1)>0\). Suppose \(U_L(-1)>0\). Then the sale is the unique maximizer of
class \(L\), and (E5) gives \(\sigma_L(-1)=1\). With \(\sigma_H(-1)=0\), \(f_H(-1)=\hat b\) and
\(f_L(-1)=\hat b+\hat\lambda\), so \(\mu(-1)=m_{\hat\lambda}\le1/2<\tau(r)\). Hence \(e(-1)=\rho\) and
\(U_L(-1)=\rho\DT m_{\hat\lambda}-k>0\). The purchase then earns
\(U_H(1)=e(1)\DT[1-\mu(1)]-k\ge\rho\DT(1-M_{\hat\lambda})-k=\rho\DT m_{\hat\lambda}-k>0\), by
Lemma~\ref{lem:a5-bounds}(a). So \(U_H(1)>0\) in both cases. The purchase is the unique maximizer of class \(H\), and
\(\sigma_H(1)=1\) under both investor rules. Lemma~\ref{lem:a5-bounds}(d) gives the three posteriors. Since
\(\tau(r)>1/2\), Lemma~\ref{lem:prices}(d) gives \(\zeta(P(0))=\zeta(P(-1))=0\) under both challenger rules.
\end{proof}

\begin{lemma}[The equilibrium set at intermediate strengths]\label{lem:a5-middle}
Let \(1/2<\tau(r)<1\), \(\rho\DT(r)<2k\) and \(\lambda_E(r)\ne\lb\). The equilibria are exactly these
tuples. Each has \(\sigma_L=\delta_0\), \(\zeta(P(0))=\zeta(P(-1))=0\), prices \(P(q)=t_0+e(q)\,w(\mu(q))\),
\(\bar\mu(P(q))=\mu(q)\), and truthful bids.
\begin{enumerate}[leftmargin=2.5em]
\item[(Z)] \(\hat\lambda=0\), \(\sigma_H=\delta_0\), \(\mu\equiv1/2\), \(\zeta\equiv0\). It always exists.
\item[(R)] \(\hat\lambda=\lambda_R(r)\), \(\sigma_H=\delta_1\),
\((\mu(-1),\mu(0),\mu(1))=\bigl(1/2,\ \hat b_0/(2\hat b_0+\hat\lambda),\ M_{\hat\lambda}\bigr)\), \(\zeta(P(1))=1\). It
exists if and only if \(\varphi_a(r)\ge0\) and \(\Lambda_r(\lb)>0\).
\item[(\(\bar{\rm R}\))] \(\hat\lambda=\lb\), with the orders, posteriors and entry of (R). It exists if and only if
\(\Lambda_r(\lb)\le0\) and \(\lambda_E(r)\le\lb\).
\item[(T)] Rule CM only: \(\hat\lambda=\lambda_E(r)\), \(\sigma_H=\delta_1\), \(\mu(1)=\tau(r)\),
\(\zeta(P(1))=(e^*-\rho)/(1-\rho)\). It exists if and only if \(\lambda_E(r)<\lb\) and \(\varphi_a(r)>0\).
\end{enumerate}
Every trade choice in these equilibria is strict, so the investor's tie rule plays no role. Entry is \(\rho\) in (Z),
\(\rho+(1-\rho)[\beta+(1/2-\beta)\hat\lambda]\) in (R) and (\(\bar{\rm R}\)), and
\(\rho+(e^*-\rho)[\beta+(1/2-\beta)\lambda_E]\) in (T).
\end{lemma}

\begin{proof}
\emph{Step 1: (Z).} (E1): \(\hat\lambda=0\) gives \(f_H=f_L=\eta\), so \(\mu\equiv1/2\). (E2): every price belief is
\(1/2\). (E3): \(\tau(r)>1/2\) gives \(B_r(1/2)<c_H\), so \(\zeta\equiv0\) with no tie. (E4): the price
\(t_0+\rho\,w(1/2)\) at every order. (E5): after a hypothetical success, \(U_H(1)=U_L(-1)=\rho\DT/2-k<0\), and the
wrong-signed orders earn at most \(-k\), so abstention is the strict best order of each class. (E6): truthful bids.
(E7): \(\Psi=0\), so \(\min\{\lb,\Psi/a\}=0=\hat\lambda\).

\emph{Step 2: (R).} Let \(\hat\lambda=\lambda_R(r)\), \(\sigma_H=\delta_1\), \(\sigma_L=\delta_0\). Then
\(f_H=(\hat b,\hat b_0,\hat b+\hat\lambda)\) and \(f_L=(\hat b,\hat b_0+\hat\lambda,\hat b)\) on \((-1,0,1)\), which give
the stated posteriors with \(\mu(0)<1/2=\mu(-1)<\mu(1)\). Suppose \(\varphi_a(r)\ge0\) and \(\Lambda_r(\lb)>0\). Then
\(\DT(1-\tau)\ge k+2a\lambda_E>k\), and \(1-\tau<1/2\) gives \(\DT>2k\), so \(\lambda_R(r)\) exists. (E1) holds by the
likelihoods. (E2): the stated entry is \(\rho\) at the posteriors \(\mu(0)\) and \(\mu(-1)\), and \(1\) at \(\mu(1)\),
so it is nondecreasing in the posterior; by the paragraph on constructed tuples in Appendix~\ref{app:beliefs}, the
three prices differ, and \(\bar\mu(P(q))=\mu(q)\). (E3): \(\Lambda_r(\lambda_E)=-[2\beta+\lambda_E(1-2\beta)]\varphi_a\le0
=\Lambda_r(\lambda_R)\), so \(\lambda_R\ge\lambda_E\) and \(\mu(1)=M_{\lambda_R}\ge\tau(r)\); rule C1 gives
\(\zeta(P(1))=1\), and rule CM forces it if \(\mu(1)>\tau\) and allows it if \(\mu(1)=\tau\); the posteriors
\(\mu(0),\mu(-1)\le1/2<\tau\) give \(\zeta=0\) there. (E4): \eqref{eq:a2-terminal}. (E5): \(U_H(1)=\DT m_{\hat\lambda}-k\),
and \(\Lambda_r(\lambda_R)=0\) gives \(\DT m_{\lambda_R}-k=2a\lambda_R>0\); \(U_L(-1)=\rho\DT\mu(-1)-k=\rho\DT/2-k<0\); the
wrong-signed orders earn at most \(-k\). So both choices are strict. (E6): truthful bids. (E7):
\(\Psi=\tfrac12(2a\lambda_R)+\tfrac12\cdot0=a\lambda_R\), and \(\Lambda_r(\lb)>0\) gives \(\lambda_R<\lb\), so
\(\min\{\lb,\Psi/a\}=\lambda_R=\hat\lambda\). Conversely, (R) needs \(\mu(1)\ge\tau\), that is \(\lambda_R\ge\lambda_E\),
that is \(\varphi_a\ge0\), and it needs \(\lambda_R<\lb\), that is \(\Lambda_r(\lb)>0\).

\emph{Step 3: (\(\bar{\rm R}\)).} The orders and posteriors are those of Step 2 with \(\hat\lambda=\lb\). If
\(\lambda_E\le\lb\), then \(\mu(1)=M_{\lb}\ge\tau\); rule C1 gives \(\zeta(P(1))=1\), and rule CM forces it if
\(\mu(1)>\tau\) and allows it if \(\mu(1)=\tau\). Then
\(\Psi=\tfrac12(\DT m_{\lb}-k)\), and \(\Lambda_r(\lb)\le0\) is \(\Psi\ge a\lb\), so \(\min\{\lb,\Psi/a\}=\lb\). The purchase earns
\(\DT m_{\lb}-k\ge2a\lb>0\), and the sale earns \(\rho\DT/2-k<0\), so the choices are strict. The other conditions
hold by direct computation: (E1) by the likelihoods of Step 2 at \(\hat\lambda=\lb\); (E2) because the stated entry is
nondecreasing in the posterior, so the three prices differ; (E4) by \eqref{eq:a2-terminal}; (E6) by truthful bids.
Both conditions are also necessary.

\emph{Step 4: (T), rule CM.} Let \(\hat\lambda=\lambda_E<\lb\), \(\sigma_H=\delta_1\), \(\sigma_L=\delta_0\) and
\(\zeta(P(1))=(e^*-\rho)/(1-\rho)\). Then \(\mu(1)=M_{\lambda_E}=\tau(r)\), so \(B_r(\mu(1))=c_H\), and rule CM
allows any \(\zeta(P(1))\). By the definition of \(e^*\), \(U_H(1)=e^*\DT(1-\tau)-k=2a\lambda_E>0\), so the purchase
is strict. The sale earns \(\rho\DT\mu(-1)-k=\rho\DT/2-k<0\), so the low class abstains strictly. Then \(\Psi=a\lambda_E\), and \(\min\{\lb,\Psi/a\}=\lambda_E\).
The preparation probability lies in \((0,1)\) if and only if \(\rho<e^*<1\). Here \(e^*<1\) is \(\varphi_a(r)>0\), and
\(e^*>\rho\) always holds, because \(\rho\DT(1-\tau)<\rho\DT/2<k<2a\lambda_E+k\). At \(e^*=1\) the tuple is (R)
with \(\lambda_R=\lambda_E\).

\emph{Step 5: no other equilibrium.} Take any equilibrium. If \(\hat\lambda=0\), then \(f_H=f_L=\eta\),
\(\mu\equiv1/2\) and \(e\equiv\rho\), so \(U_H(1)=U_L(-1)=\rho\DT/2-k<0\), and (E5) gives
\(\sigma_H=\sigma_L=\delta_0\): this is (Z). If \(\hat\lambda>0\), Lemma~\ref{lem:a5-buy} gives \(\sigma_H=\delta_1\),
\(\mu(1)=M_{\hat\lambda}\), \(\mu(-1)\le1/2\) and \(e(-1)=\rho\). Then \(U_L(-1)\le\rho\DT/2-k<0\), so
\(\sigma_L=\delta_0\), and \(\Psi=\tfrac12\bigl(e(1)\DT m_{\hat\lambda}-k\bigr)^+\). Three cases remain.
\begin{itemize}[leftmargin=2em]
\item \(\hat\lambda<\lambda_E\): \(\mu(1)<\tau\), so \(e(1)=\rho\), and \(\rho\DT m_{\hat\lambda}<\rho\DT/2<k\), so
\(\Psi=0\). Lemma~\ref{lem:a5-joint} then gives \(\hat\lambda=0\), a contradiction.
\item \(\hat\lambda>\lambda_E\): \(\mu(1)>\tau\), so \(e(1)=1\) under both rules. If \(\hat\lambda<\lb\), then
\(a\hat\lambda=\Psi>0\), which is \(\Lambda_r(\hat\lambda)=0\), so \(\hat\lambda=\lambda_R\): this is (R). If
\(\hat\lambda=\lb\), then \(\Psi\ge a\lb\), which is \(\Lambda_r(\lb)\le0\): this is (\(\bar{\rm R}\)).
\item \(\hat\lambda=\lambda_E\): under rule C1, \(e(1)=1\), and the previous case applies with \(\lambda_R=\lambda_E\).
Under rule CM, \(e(1)=e\in[\rho,1]\). Since \(\lambda_E\ne\lb\), research is interior, and
\(a\lambda_E=\tfrac12\bigl(e\DT(1-\tau)-k\bigr)^+\), so \(e=e^*\): this is (T), or (R) when \(e^*=1\).
\end{itemize}
In each listed equilibrium with positive research, the purchase earns \(2a\hat\lambda>0\) or more.
The sale earns \(\rho\DT/2-k<0\). Step 1 gives strict abstention in (Z), so every trade choice is strict.
A purchase has probability
\(\tfrac12(\hat b+\hat\lambda)+\tfrac12\hat b=\beta+(1/2-\beta)\hat\lambda\), and only the purchase can raise entry
above \(\rho\). This gives the three entry formulas.
\end{proof}

\begin{lemma}[Strong incumbents]\label{lem:a5-strong}
Let \(\tau(r)>1/2\), \(M_{\lb}\ne\tau(r)\) and
\begin{equation*}\tag{U}\label{eq:a5-U}
\rho\DT(r)\,m_{\lb}-k>a\lb .
\end{equation*}
Under every pair of tie rules the equilibrium is unique: \(\hat\lambda=\lb\), \(\sigma_H=\delta_1\),
\(\sigma_L=\delta_{-1}\), \((\mu(-1),\mu(0),\mu(1))=(m_{\lb},1/2,M_{\lb})\), \(\zeta(P(0))=\zeta(P(-1))=0\), and
\(\zeta(P(1))=\1\{M_{\lb}>\tau(r)\}\), with \(P(q)=t_0+e(q)\,w(\mu(q))\), \(\bar\mu(P(q))=\mu(q)\) and truthful bids.
Entry is \(\rho+(1-\rho)(\bar b+\lb/2)\), with \(\bar b=(1-\lb)\beta\), if \(M_{\lb}>\tau(r)\), and \(\rho\) if
\(M_{\lb}<\tau(r)\). This is the equilibrium of Proposition~\ref{prop:ccc} with presence \(\lb\).
\end{lemma}

\begin{proof}
\emph{Uniqueness.} Take any candidate. By Lemma~\ref{lem:a5-bounds}(a), every posterior lies in
\([m_{\lb},M_{\lb}]\) whatever the orders, and \(e\ge\rho\) after every price. Hence
\(U_H(1)\ge\rho\DT(1-M_{\lb})-k\) and \(U_L(-1)\ge\rho\DT m_{\lb}-k\). Both lower bounds equal
\(\rho\DT m_{\lb}-k>a\lb\ge0\) by \eqref{eq:a5-U}. So the purchase and the sale are the strict best orders after a
success, under both investor rules. Then \(\Psi=\tfrac12[U_H(1)+U_L(-1)]>a\lb\), and Lemma~\ref{lem:a5-joint} gives
\(\hat\lambda=\lb\). With full orders, \(f_H=(\bar b,\bar b_0,\bar b+\lb)\) and \(f_L=(\bar b+\lb,\bar b_0,\bar b)\), so the
posteriors are \((m_{\lb},1/2,M_{\lb})\). Since \(\tau(r)>1/2\), \(\zeta(P(0))=\zeta(P(-1))=0\). Since
\(M_{\lb}\ne\tau(r)\), \(\zeta(P(1))\) is \(1\) or \(0\) with no tie, under both challenger rules. Every component of
the tuple is fixed.

\emph{Verification.} (E1): the likelihoods above. (E2): the stated entry is nondecreasing in the posterior, so the
three prices differ, and each price reveals its order. (E3): \(B_r(M_{\lb})\) lies strictly on one side of \(c_H\),
and \(B_r(1/2),B_r(m_{\lb})<c_H\). (E4): \eqref{eq:a2-terminal}. (E5): the strict bounds above, read at the tuple.
(E6): truthful bids. (E7): \(\Psi>a\lb\). The purchase has probability \(\tfrac12(\bar b+\lb)+\tfrac12\bar b=\bar b+\lb/2\),
and it recruits exactly when \(M_{\lb}>\tau(r)\).
\end{proof}

\begin{proof}[Proof of Proposition~\ref{prop:research}(b)]
Throughout, \(a=\reg{res_a}\) and \(\lb=\reg{res_cap}\).

\emph{Part (b)(i).} Take any candidate at \(r_0\). The gross return of a nonzero order after a success is at most
\(e(q)\DT\max\{\mu(q),1-\mu(q)\}\le\DT\), and \(\DT(r_0)=\reg{res_Delta_r0}<k\), with margin
\(\reg{margin_res_weak_B}\). So \(U_j(q)<0\) for \(q\ne0\) against every price function and entry rule. Then
\(\Psi=0\), Lemma~\ref{lem:a5-joint} gives \(\hat\lambda=0\), and (E5) gives \(\sigma_H=\sigma_L=\delta_0\). Every
posterior is \(1/2\), and \(B_{r_0}(1/2)<c_H\), with margin \(\reg{margin_res_half_B}\), so \(\zeta\equiv0\) and entry
is \(\rho\). Step 1 of the proof of Lemma~\ref{lem:a5-middle} verifies this tuple. The general statement behind this
part needs both \(\DT(r)<k\) and \(B_r(1/2)<c_H\): with \(c_H\) small enough that \(B_r(1/2)>c_H\), zero research is
still unique, but both cost types prepare and entry is one.

\emph{Part (b)(ii).} At each \(r\in\{\reg{str_r155},\reg{str_r160},\reg{str_r165}\}\), the hypotheses of
Lemma~\ref{lem:a5-middle} hold: \(1/2<\tau(r)<1\); \(k-\rho\DT(r)/2\) equals \(\reg{res_ULgap_r155}\),
\(\reg{res_ULgap_r160}\) and \(\reg{res_ULgap_r165}\); and \(\lambda_E(r)<\lb\). Also \(\lambda_R(r)-\lambda_E(r)\)
equals \(\reg{res_lamgap_r155}\), \(\reg{res_lamgap_r160}\) and \(\reg{res_lamgap_r165}\), so \(\varphi_a(r)>0\); and
\(\lb-\lambda_R(r)\) equals \(\reg{res_capgap_r155}\), \(\reg{res_capgap_r160}\) and \(\reg{res_capgap_r165}\), so
\(\Lambda_r(\lb)>0\). Lemma~\ref{lem:a5-middle} therefore gives the equilibrium set \(\{{\rm Z},{\rm R}\}\) under rule C1,
and \(\{{\rm Z},{\rm R},{\rm T}\}\) under rule CM, under either investor rule. Table~\ref{tab:research} gives the
values. The research values come from bisection of \(\Lambda_r\) with exact signs; \(\DT\), \(\tau\), \(\lambda_E\) and
\(e^*\) are exact rationals.

\emph{Part (b)(iii).} At calibration B, \(M_{\lb}=\eu/(1+\eu)\), and \(\tau(r_1)<M_{\lb}<\tau(r_2)\), with margins
\(\reg{margin_res_tau_r1_B}\) and \(\reg{margin_res_tau_r2_B}\) in posterior. Condition~\eqref{eq:a5-U} is
\(\DT(r)>(k+a\lb)(1+\eu)/\rho\), which holds for every \(r>r_U=\fr\bigl((k+a\lb)(1+\eu)/\rho\bigr)=\reg{res_rUR}\);
its margin is \(\reg{margin_res_bound_r1_B}\) at \(r_1\) and \(\reg{margin_res_bound_r2_B}\) at \(r_2\).
Lemma~\ref{lem:a5-strong} gives the unique equilibrium at both strengths, with entry \(\reg{res_entry_r1}\) at
\(r_1\) and \(\reg{res_entry_r2}\) at \(r_2\).

Every trade choice at the six strengths is strict, by Step 5 of the proof of Lemma~\ref{lem:a5-middle} and by the
bounds of Lemma~\ref{lem:a5-strong}. The investor's tie rule therefore plays no role.
\end{proof}

\begin{proof}[Proof of Proposition~\ref{prop:research}(a)]
The polynomial \(\Lambda_r(\lambda)\) is smooth in \((\lambda,r)\), since its coefficients are affine in \(\DT(r)\). For
\(\lambda\ge0\),
\[
\partial_\lambda \Lambda_r=4a(1-2\beta)\lambda+4a\beta+k(1-2\beta)+\beta\DT>0,\qquad
\partial_r\Lambda_r=-\beta(1-\lambda)\DT'(r)<0,
\]
because \(\lambda<1\) and \(\DT'(r)=\tfrac12(1-\ell^2/r^2)>0\). On \(\{r:\DT(r)>2k\}\), the implicit function theorem
gives
\[
\lambda_R'(r)=\frac{\beta\,[1-\lambda_R(r)]\,\DT'(r)}{\partial_\lambda \Lambda_r(\lambda_R(r))}>0 .
\]
On the branch (R), Lemma~\ref{lem:a5-middle} gives \(\Ent_R=\rho+(1-\rho)[\beta+(1/2-\beta)\lambda_R]\). Since
\(\beta<1/2\), \(\Ent_R\) rises strictly with \(\lambda_R\) and hence with \(r\):
\(\Ent_R'=(1-\rho)(1/2-\beta)\lambda_R'>0\). Class \(H\) buys with probability
\(\hat b+\lambda_R=\beta+(1-\beta)\lambda_R\), and \(W_H=1\) because \(h>r\). So ownership is
\(\tfrac12[\rho+(1-\rho)(\beta+(1-\beta)\lambda_R)]\), which also rises strictly. These signs hold for every
\(0<\beta\le1/3\), wherever the branch is an equilibrium.
\end{proof}

\paragraph{The branch at calibration B.}
At \(a=\reg{res_a}\), (R) exists exactly on \([r_R,r_{\lb})=[\reg{res_rR},\reg{res_rcap})\). The proof has four
steps.
\begin{enumerate}[leftmargin=2em]
\item For \(r\le r_0\), \(\DT(r)\le\DT(r_0)<k\), so part (b)(i) applies, and (R) does not exist.
\item The margin \(\varphi_{a}\) rises on \([r_0,7/4]\). By the formula for \(d\lambda_E/d\tau\),
\(\varphi_a'=\DT'(1-\tau)-\DT\tau'-2a\beta\tau'/\{\beta+(1-2\beta)(1-\tau)\}^2\), with
\(\tau'(r)=3(8r^2-r+2)/[2(20r-r^2-1)^2]>0\). On this interval \(\DT'\), \(\DT\) and \(\tau\) rise, so
\(\DT'\ge\reg{res_bound_dDelta}\), \(1-\tau\ge1-\tau(7/4)=\reg{res_bound_tau74}\), \(\DT\le\DT(7/4)=\reg{res_bound_Delta74}\)
and \(\tau'\le\reg{res_bound_taup}\). These give \((\DT(1-\tau))'\ge\reg{res_dphi_lead}\) and, at
\(a=\reg{res_a}\), \(\varphi_a'\ge\reg{res_phip_lb}>0\). The bound stays positive for every
\(a<a_{\rm mono}=\reg{res_amono}\). Exact rational evaluation gives \(\varphi_a(r_0)<0<\varphi_a(7/4)\). So
\(\varphi_a\) has the unique root \(r_R\) on \([r_0,7/4]\), and bisection with exact signs encloses it.
\item For \(r_R\le r<r_{\lb}\): \(r<r_N\), so \(\rho\DT<2k\), and \(1/2<\tau<1\); also \(\varphi_a(r)\ge0\). Since
\(m_{\lb}=1/(1+\eu)\), \(\Lambda_r(\lb)=0\) is \(\DT(r)=(2a\lb+k)(1+\eu)\), and \(r_{\lb}=\fr\bigl((2a\lb+k)(1+\eu)\bigr)\).
Since \(\lambda_R\) rises and \(\lambda_R(r_{\lb})=\lb\), \(\Lambda_r(\lb)>0\) for \(r<r_{\lb}\), and (R) exists by
Lemma~\ref{lem:a5-middle}. On \((r_0,r_R)\), \(\varphi_a<0\), and (R) does not exist.
\item For \(r\ge r_{\lb}\), take any equilibrium with \(e(1)=1\) and \(\hat\lambda<\lb\). Lemma~\ref{lem:a5-buy}
gives \(\sigma_H=\delta_1\) and \(\mu(1)=M_{\hat\lambda}\). The low class earns at least zero, so
\(a\hat\lambda=\Psi\ge\tfrac12(\DT m_{\hat\lambda}-k)\), that is \(\Lambda_r(\hat\lambda)\ge0\). Hence
\(\hat\lambda\ge\lambda_R(r)\ge\lambda_R(r_{\lb})=\lb\), a contradiction. This step also covers \(r>r_N\), where the
low class may sell.
\end{enumerate}
Zero research is an equilibrium exactly when \(\rho\DT\le2k\), that is \(r\le r_N\), and
\(r_{\lb}<r_N=\reg{res_rN}\), with margin \(\reg{margin_res_rN_B}\). Along the branch, research rises from
\(\lambda_E(r_R)=\reg{res_lambda_rR}\) to the cap, and entry rises from \(\reg{res_entry_rR}\) to
\(\reg{res_entry_cap}\). On \([r_{\lb},r_N)\), \(\Lambda_r(\lb)\le0\) and \(\lambda_E(r)\le\lambda_E(7/4)=\reg{res_lamE_r175}<\lb\),
so (\(\bar{\rm R}\)) exists by Lemma~\ref{lem:a5-middle}, with entry \(\reg{res_entry_cap}\). Under rule CM, (T) exists on
\((r_R,r_N)\), since \(\varphi_a>0\) and \(\lambda_E<\lb\) there.

\begin{proposition}[The research parameter region]\label{prop:researchregion}
Fix \(\beta=\reg{beta_B}\) and the other primitives of \(B\), and vary the cap \(\lb\) and the cost \(a\). The equilibrium patterns of
Proposition~\ref{prop:research}(b), with \(\lambda_E(r)<\lambda_R(r)<\lb\) at each intermediate strength, hold with
strict margins if and only if
\[
\lambda_E(r_1)<\lb<\lambda_E(r_2)\qquad\text{and}\qquad a_{\mathrm{lo}}(\lb)<a<a_{\mathrm{up}},
\]
where
\[
a_{\mathrm{lo}}(\lb)=\frac{\DT(\reg{str_r165})\,m_{\lb}-k}{2\lb},\qquad
a_{\mathrm{up}}=\frac{\DT(\reg{str_r155})\,[1-\tau(\reg{str_r155})]-k}{2\lambda_E(\reg{str_r155})}.
\]
The cap strip is \((\reg{res_cap_lo},\reg{res_cap_hi})\) to six places. At the cap \(\lb=\reg{res_cap}\), the cost
window is \((\reg{res_a_lo},\reg{res_a_up})\); at \(a=\reg{res_a}\), the cap range is
\((\reg{res_cap_lo_a},\reg{res_cap_hi})\). Entry at \(r_1\) is \(\rho+(1-\rho)[\beta+(1/2-\beta)\lb]\).
With \(1+u=\eu\), the cap \(\lb=u\beta/(1+u\beta)\) and \(a=\reg{res_a}\),
these patterns hold for every \(\beta\in(\reg{res_beta_lo},1/3]\).
\end{proposition}

The lower cost bound keeps research below the cap at \(\reg{str_r165}\). The upper cost bound keeps recruitment after a
purchase at \(\reg{str_r155}\). The cap strip gives strict recruitment at \(r_1\) and strict exclusion at \(r_2\).

\begin{proof}[Proof of Proposition~\ref{prop:researchregion}]
Fix \(\beta=\reg{beta_B}\) and the other primitives of calibration B. The required equilibrium patterns are those of Proposition~\ref{prop:research}(b)
with \(\lambda_E(r)<\lambda_R(r)<\lb\) at each intermediate strength.

\emph{Sufficiency.} Let the two conditions hold. Part (b)(i) holds for every \((\lb,a)\), since \(\DT(r_0)<k\). For part
(b)(ii), the hypotheses \(1/2<\tau<1\) and \(\rho\DT<2k\) of Lemma~\ref{lem:a5-middle} do not involve \((\lb,a)\) and hold
at the three strengths. The condition \(a<a_{\rm up}\) is \(\varphi_a(\reg{str_r155})>0\). Since
\(a_{\rm up}<a_{\rm mono}\), \(\varphi_a\) rises on \([r_0,7/4]\), so \(\varphi_a>0\) at \(\reg{str_r160}\) and
\(\reg{str_r165}\) too. The condition \(a>a_{\rm lo}(\lb)\) is \(\Lambda_{\reg{str_r165}}(\lb)>0\), that is
\(\lambda_R(\reg{str_r165})<\lb\). Since \(\lambda_R\) rises in \(r\), \(\lambda_R<\lb\) at the two lower strengths
too. Then \(\lambda_E<\lambda_R<\lb\), and Lemma~\ref{lem:a5-middle} gives the sets. For part (b)(iii), \(M_\lambda\) rises
in \(\lambda\) and \(M_{\lambda_E(r)}=\tau(r)\), so the strip gives \(\tau(r_1)<M_{\lb}<\tau(r_2)\). Condition
\eqref{eq:a5-U} at \(r_1\) is \(a<a_{U}(r_1;\lb):=[\rho\DT(r_1)m_{\lb}-k]/\lb\), and at \(r_2\) it is
\(a<a_{U}(r_2;\lb):=[\rho\DT(r_2)m_{\lb}-k]/\lb\). Both bounds fall in \(\lb\), so on the strip they exceed their
values at \(\lb=\lambda_E(r_2)\), which are \(\reg{res_region_aU3_6}\) and \(\reg{res_region_a36_6}\). Both exceed
\(a_{\rm up}=\reg{res_a_up}\), so Lemma~\ref{lem:a5-strong} applies at both strengths.
Its entry formula gives \(\rho+(1-\rho)[\beta+(1/2-\beta)\lb]\) at \(r_1\).

\emph{Necessity.} Each failure of a condition breaks one claim.
\begin{itemize}[leftmargin=2em]
\item \(\lb\le\lambda_E(r_1)\): every posterior at \(r_1\) is at most \(M_{\lb}\le\tau(r_1)\), so entry above \(\rho\)
needs a challenger tie, and part (b)(iii) fails with strict margins.
\item \(\lb\ge\lambda_E(r_2)\): research at the cap gives \(\mu(1)=M_{\lb}\ge\tau(r_2)\), so rule C1 raises entry above
\(\rho\) at \(r_2\), and part (b)(iii) fails.
\item \(a\ge a_{\rm up}\): \(\varphi_a(\reg{str_r155})\le0\), so (R) is absent at \(\reg{str_r155}\), or has
\(\mu(1)=\tau\) with a challenger tie.
\item \(a\le a_{\rm lo}(\lb)\): \(\Lambda_{\reg{str_r165}}(\lb)\le0\), so research at \(\reg{str_r165}\) sits at the cap.
\end{itemize}
The cost window is nonempty on the whole strip: \(a_{\rm lo}\) falls in \(\lb\), because its numerator falls and stays
positive there, and \(a_{\rm lo}(\lambda_E(r_1))=\reg{res_region_alo_0}<a_{\rm up}\). The cap strip is
\((\reg{res_cap_lo},\reg{res_cap_hi})\); the margins of \(\lb=\reg{res_cap}\) are \(\reg{margin_res_caplo_B}\) and
\(\reg{margin_res_caphi_B}\), and the margins of \(a=\reg{res_a}\) in its window are \(\reg{margin_res_alo_B}\) and
\(\reg{margin_res_aup_B}\). The upper cap margin corresponds to the posterior margin \(\tau(r_2)-M_{\lb}=
\reg{margin_res_tau_r2_B}\); the two numbers measure one boundary in two coordinates.

\emph{General \(\beta\).} Set \(\lb=\lambda_B(\beta)=u\beta/(1+u\beta)\) with \(1+u=\eu\), and keep \(a=\reg{res_a}\).
Then \(u_{\lb}=u\) for every \(\beta\), so \(M_{\lb}=\eu/(1+\eu)\) and \(m_{\lb}=1/(1+\eu)\) do not depend on \(\beta\),
and the endpoint conditions \(\tau(r_1)<M_{\lb}<\tau(r_2)\) hold for every \(\beta\). Condition \eqref{eq:a5-U} at
\(r_1\) holds for every \(\beta\le1/3\), because \(\lambda_B(\beta)\) rises in \(\beta\) and
\(\rho\DT(r_1)/(1+\eu)-k-a\lambda_B(1/3)>0\). At the intermediate strengths,
\(a_{\rm lo}=[\DT(\reg{str_r165})/(1+\eu)-k]/[2\lambda_B(\beta)]\) falls in \(\beta\), so \(a_{\rm lo}<a\) exactly when
\(\beta>\beta_{\rm lo}=\reg{res_beta_lo}\). The upper bound \(a_{\rm up}(\beta)\) falls in \(\beta\), because
\(\partial\lambda_E/\partial\beta=(2\tau-1)(1-\tau)/\{\beta+(1-2\beta)(1-\tau)\}^2>0\), and it exceeds \(a\) for every
\(\beta<\reg{res_beta_hi}\). The monotonicity bound \(a_{\rm mono}(\beta)\) falls in \(\beta\) on \((0,1/3]\) and equals
\(\reg{res_amono_third}>a\) at \(\beta=1/3\). The steps above then apply at each \(\beta\in(\reg{res_beta_lo},1/3]\).
Strict growth of entry across the three strengths alone holds on a wider set, because research can reach the cap at
\(\reg{str_r165}\) while entry still rises.
\end{proof}

% STRUCTURE: Remark 2 <- Section 5, with the threshold-equilibrium entry of Table 4 (rule CM)
\begin{remark}[Tie rules with research]\label{rem:researchties}
Research never has a tie, and Proposition~\ref{prop:research}(b) needs no investor tie rule. Under the challenger rule CM,
each intermediate strength has one more equilibrium. In it, research is \(\lambda_E(r)\), and the purchase posterior
equals the threshold. The indifferent challenger prepares with the probability that makes this research optimal. Its
entry falls across the three strengths: it is \(\reg{res_th_entry_r155}\), \(\reg{res_th_entry_r160}\) and
\(\reg{res_th_entry_r165}\) at \(r=\reg{str_r155}\), \(\reg{str_r160}\) and \(\reg{str_r165}\). Above \(r_U\), the whole tuple is unique under C1,
and under CM at every strength except the ceiling \(r_C=\reg{res_rC}\). At \(r_C\) under CM, entry can take any value
in \([\rho,\reg{res_entry_cap}]\).
\end{remark}

\begin{proof}[Proof of Remark~\ref{rem:researchties}]
Every trade choice is strict at the six strengths (proof of Proposition~\ref{prop:research}(b)), and research has no
tie (Lemma~\ref{lem:a5-joint}). The threshold equilibrium is (T) of Lemma~\ref{lem:a5-middle}; its entry is
\(\reg{res_th_entry_r155}\), \(\reg{res_th_entry_r160}\) and \(\reg{res_th_entry_r165}\) at the three strengths, with
\(e^*\) equal to \(\reg{res_th_estar_r155}\), \(\reg{res_th_estar_r160}\) and \(\reg{res_th_estar_r165}\). For
\(r>r_U\), the bounds of Lemma~\ref{lem:a5-strong} force full orders and research at the cap in every candidate.
Under rule C1, or under rule CM with \(M_{\lb}\ne\tau(r)\), the whole tuple is then unique. At \(r=r_C\),
\(M_{\lb}=\tau(r_C)\), the high-cost challenger is indifferent after a purchase, and rule CM admits every
\(\zeta(P(1))\in[0,1]\), with \(\zeta(P(0))=\zeta(P(-1))=0\). Each value has its own competitive purchase price, and
entry \(\rho+(1-\rho)\zeta(P(1))(\bar b+\lb/2)\) fills \([\rho,\reg{res_entry_cap}]\). Research stays at the cap.
\end{proof}

\paragraph{Access to prices with research.}
At \(r_1\), compare the economy above with one in which the challenger does not see the price. In that economy the
challenger's belief stays at \(1/2\), and \(B_{r_1}(1/2)<c_H\) gives entry \(\rho\) after every order. The bounds of
Lemma~\ref{lem:a5-strong} hold with \(e\equiv\rho\), so \eqref{eq:a5-U} again forces full orders and research at the
cap, uniquely. The order law, the research cost \((a/2)\lb^2\) and the trade cost \(k\lb\) are therefore the same in
both economies. Entry differs only after a purchase, where it is \(1\) with access and \(\rho\) without. The purchase
has probability \(\tfrac12(\bar b+\lb)\) in class \(H\) and \(\tfrac12\bar b\) in class \(L\), so access raises
expected proceeds by \((1-\rho)(\bar b+\lb/2)[t_L-t_0+\DT M_{\lb}]=\reg{res_dR_r1}\). For \(\theta\ge p\),
\(\max\{\theta,R\}-R\1\{R\ge p\}=(\theta-\max\{p,R\})_++p\1\{R<p\}\), so a prepared challenger of class \(j\) adds
\(g_j+p\Pr(R<p)=g_j+p^2/r_1\) to allocation value. Net of the added preparation cost, access raises surplus by
\((1-\rho)(\bar b+\lb/2)[B_{r_1}(M_{\lb})-c_H+p^2/r_1]=\reg{res_dW_r1}\). Both are positive, because
\(B_{r_1}(M_{\lb})>c_H\).

\paragraph{Certificates.}
The exact program is \nolinkurl{P3_research/P3_research.py} (run log
\nolinkurl{P3_research/P3_research.log}). It uses rational arithmetic for
every rational quantity and rational intervals for every quantity that depends on \(\eu\). It builds the tuple of
every equilibrium of Lemmas~\ref{lem:a5-middle} and~\ref{lem:a5-strong} for each pair of tie rules and recomputes
(E1) to (E7) at the tuple. It rounds every display outward from exact rational endpoints and writes the endpoints to
\nolinkurl{P3_research/P3_enclosures.json}. The independent program
\nolinkurl{XP3/recompute.py} (run log
\nolinkurl{XP3/recompute.log}) computes the research roots from the
quadratic formula with rational bounds on every square root. It bounds \(\eu\) between \(\sum_{n\le60}1/n!\) and that
sum plus \(62/(61\cdot61!)\), and it records exact rational endpoints for every printed value in
\nolinkurl{XP3/exact_enclosures.json}. These rational endpoints are the
certificate of each display. The scan
\nolinkurl{P3_research/P3_float_scan.py} (run log
\nolinkurl{P3_research/P3_float_scan.log}) evaluates a grid of research
intensities and order mixtures at each of the six strengths without the classification of
Lemma~\ref{lem:a5-middle}. It finds exactly the equilibria of the lemmas under rule C1. This scan enters no proof.
